\SGAchapter{VIII}{忠实平坦下降}{Descente fidèlement plate}
\SourcePage{169}
\SGAsection{1}{拟凝聚模的下降}
设 $\operatorname{Sch}$ 为\Term{prescheme}{预概形}\Term{category}{范畴}。按照第 VI 讲 11.b 的方式，可以把由\Term{prescheme}{预概形} $X$ 及其上的模 $F$ 所组成的对 $(X,F)$ 的\Term{category}{范畴}，看成 $\operatorname{Sch}$ 上的一个\Term{fibred-category}{纤维范畴}；其中\Term{morphism}{态射}如该处一样，借助环空间\Term{morphism}{态射}下模的\Term{direct-image}{正像}来定义。对于 $\operatorname{Sch}$ 中的\Term{morphism}{态射} $f:X\to Y$，相应的\Term{base-change}{基变换}\Term{functor}{函子}就是模在 $f$ 下的\Term{inverse-image}{逆像}\Term{functor}{函子}。（注意，上述\Term{fibred-category}{纤维范畴}在 $X\in\operatorname{Ob}(\operatorname{Sch})$ 处的\Term{fibre-category}{纤维范畴}，是 $X$ 上模\Term{category}{范畴}的\Term{opposite-category}{反范畴}。）由于\Term{quasi-coherent-module}{拟凝聚模}的\Term{inverse-image}{逆像}仍然\Term{quasi-coherent}{拟凝聚}，由所有满足 $F$ \Term{quasi-coherent}{拟凝聚}的对 $(X,F)$ 组成的\Term{full-subcategory}{满子范畴}，是上述\Term{fibred-category}{纤维范畴}的一个\Term{fibre}{纤维}子\Term{category}{范畴}。（反之，若不对 $f$ 作假设，\Term{quasi-coherent-module}{拟凝聚模}的\Term{direct-image}{正像}一般并不\Term{quasi-coherent}{拟凝聚}。）我们就把这个\Term{fibred-category}{纤维范畴}称为\Term{prescheme}{预概形}上的\Term{quasi-coherent-module}{拟凝聚模}\Term{fibred-category}{纤维范畴}。

还需回顾：环空间\Term{morphism}{态射} $f:X\to Y$ 若为\Term{flat}{平坦}的，即对每个 $x\in X$，$\mathcal O_{X,x}$ 都是 $\mathcal O_{Y,f(x)}$ 上的\Term{flat}{平坦}模（第 IV 讲），而且 $f$ 为\Term{surjectivity}{满射}，就称 $f$ \Term{faithfully-flat}{忠实平坦}。若每个\Term{quasi-compact}{拟紧}子集在 $f$ 下的\Term{inverse-image}{逆像}都\Term{quasi-compact}{拟紧}，就称 $f$ 为\Term{quasi-compact}{拟紧}\Term{morphism}{态射}；当 $f$ 是\Term{prescheme}{预概形}\Term{morphism}{态射}时，这也等价于：$Y$ 中每个\Term{affine-open}{仿射开集}的\Term{inverse-image}{逆像}，都是 $X$ 中\Term{finite}{有限}多个\Term{affine-open}{仿射开集}的并。

\begin{SGAresult}{定理}{1.1}
设 $\mathcal F$ 为\Term{prescheme}{预概形}上\Term{quasi-coherent-module}{拟凝聚模}的\Term{fibred-category}{纤维范畴}。设\Term{prescheme}{预概形}\Term{morphism}{态射} $g:S'\to S$ \Term{faithfully-flat}{忠实平坦}且\Term{quasi-compact}{拟紧}。则 $g$ 是一个 $\mathcal F$-\Term{effective-descent-morphism}{有效下降态射}。
\end{SGAresult}
回顾一下，\footnote{这里采用\Term{descent}{下降}的一般理论；J. Giraud 的文章对此有详细论述，见《油印版说明》脚注 (5)，下文将这篇文章记为 [D]。简要论述另见 [2]。}这意味着以下两件事。

\SourcePage{170}
\begin{SGAresult}{推论}{1.2（模\Term{homomorphism}{同态}的\Term{descent}{下降}）}
设 $g:S'\to S$ 是\Term{faithfully-flat}{忠实平坦}且\Term{quasi-compact}{拟紧}的\Term{prescheme}{预概形}\Term{morphism}{态射}，$F,G$ 为 $S$ 上的两个\Term{quasi-coherent-module}{拟凝聚模}，$F',G'$ 为它们在 $S'$ 上的\Term{inverse-image}{逆像}，$F'',G''$ 为它们在 $S''=S'\times_S S'$ 上的\Term{inverse-image}{逆像}。由 $g,p_1,p_2$ 下的\Term{base-change}{基变换}\Term{functor}{函子}确定下列集合映射图，其中 $p_1,p_2:S'\times_S S'\rightrightarrows S'$ 是两个投影：
\[
\operatorname{Hom}_S(F,G)\longrightarrow
\operatorname{Hom}_{S'}(F',G')\rightrightarrows
\operatorname{Hom}_{S''}(F'',G'').
\]
该图是\Term{exact}{正合}的，即第一集合与第二集合中在两条通往第三集合的映射下取值相同的元素所成的集合之间，给出了一个\Term{bijection}{双射}。
\end{SGAresult}
换言之，沿 $g$ 的\Term{base-change}{基变换} $F\mapsto F'$，给出一个\Term{fully-faithful}{全忠实函子}：从 $S$ 上\Term{quasi-coherent-module}{拟凝聚模}的\Term{category}{范畴}，到带有相对于 $g$ 的\Term{descent-datum}{下降数据}的 $S'$ 上\Term{quasi-coherent-module}{拟凝聚模}的\Term{category}{范畴}。此外还有：

\begin{SGAresult}{推论}{1.3（模的\Term{descent}{下降}）}
对 $S'$ 上每个\Term{quasi-coherent-module}{拟凝聚模} $F'$，其上相对于 $g$ 的每个\Term{descent-datum}{下降数据}都是有效的。也就是说，连同其\Term{descent-datum}{下降数据}，$F'$ \Term{isomorphism}{同构}于 $S$ 上某个\Term{quasi-coherent-module}{拟凝聚模}沿 $g$ 的\Term{inverse-image}{逆像}；由 1.2，这个模在唯一\Term{isomorphism}{同构}的意义下确定。
\end{SGAresult}
换言之，前述\Term{fully-faithful}{全忠实函子}实际上是一个\Term{equivalence-of-categories}{范畴等价}。实际使用时，这意味着给定 $S$ 上一个\Term{quasi-coherent-module}{拟凝聚模}，等价于给定 $S'$ 上一个\Term{quasi-coherent-module}{拟凝聚模}，并赋予相对于 $g$ 的\Term{descent-datum}{下降数据}。

\textbf{1.1 的证明。}先设 $T$ 为一个 $S$-\Term{prescheme}{预概形}，它作为 $S$-\Term{prescheme}{预概形}\Term{isomorphism}{同构}于一族\Term{covering}{覆盖} $S$ 的开子\Term{prescheme}{预概形} $S_i$ 的不交并。显然，结构\Term{morphism}{态射} $T\to S$ 是 $\mathcal F$-\Term{effective-descent-morphism}{有效下降态射}：这恰好是说，给定 $S$ 上一个\Term{quasi-coherent-module}{拟凝聚模} $F$，等价于给定每个 $S_i$ 上的\Term{quasi-coherent-module}{拟凝聚模} $F_i$，以及满足熟知的上\Term{chain-condition}{链条件}的\Term{gluing-isomorphism}{粘合同构}
\[
\varphi_{ji}:F_i|_{S_i\cap S_j}\longrightarrow F_j|_{S_i\cap S_j}.
\]
由第 VII 讲第 8 节，要检验 $g:S'\to S$ 是 $\mathcal F$-\Term{effective-descent-morphism}{有效下降态射}，只需对由\Term{base-change}{基变换} $T\to S$ 得到的\Term{morphism}{态射} $g_T:T'=T\times_S S'\to T$ 检验这一点。（注意，对 $T\to S$ 的上述假设在任意\Term{base-change}{基变换}下保持，故 $T\to S$ 实际上是一个普遍的 $\mathcal F$-\Term{effective-descent-morphism}{有效下降态射}。）取一族\Term{covering}{覆盖} $S$ 的\Term{affine-open}{仿射开集}作为 $S_i$，便归结为 $S$ \Term{affine}{仿射}的情形。

此时 $S'$ 为\Term{finite}{有限}多个\Term{affine-open}{仿射开集}的并。取它们作为 $S$-\Term{scheme}{概形}的不交并，就得到一个\Term{affine}{仿射} $S$-\Term{scheme}{概形} $S_1$，以及\Term{flat}{平坦}且\Term{surjectivity}{满射}的 $S$-\Term{morphism}{态射} $S_1\to S'$。因此 $S_1$ 在 $S$ 上也\Term{faithfully-flat}{忠实平坦}。若能证明\Term{faithfully-flat}{忠实平坦}的\Term{affine-morphism}{仿射态射}为 $\mathcal F$-\Term{effective-descent-morphism}{有效下降态射}，从而为普遍严格的 $\mathcal F$-\Term{descent-morphism}{下降态射}——因为这一假设确实在\Term{base-change}{基变换}下保持——
\SourcePage{171}
就特别可知结构\Term{morphism}{态射} $S_1\to S$ 是普遍严格的 $\mathcal F$-\Term{descent-morphism}{下降态射}。又因存在 $S$-\Term{morphism}{态射} $S_1\to S'$，由 [D] 可知 $g:S'\to S$ 是严格的 $\mathcal F$-\Term{descent-morphism}{下降态射}。

这样就归结为 $g$ \Term{affine}{仿射}的情形；如前所见，此时还可假设 $S$ \Term{affine}{仿射}，故可假设 $S,S'$ 都\Term{affine}{仿射}。在这种情况下，1.2 等价于以下引理。

\begin{SGAresult}{引理}{1.4}
设 $A$ 为环，$A'$ 为\Term{faithfully-flat}{忠实平坦}的 $A$-代数，$M,N$ 为两个 $A$-模。以 $M',N'$ 表示沿 $A\to A'$ 作环变换得到的 $A'$-模，以 $M'',N''$ 表示沿 $A\to A''$ 作环变换得到的 $A''=A'\otimes_A A'$-模。则集合映射序列
\[
\operatorname{Hom}_A(M,N)\longrightarrow
\operatorname{Hom}_{A'}(M',N')\rightrightarrows
\operatorname{Hom}_{A''}(M'',N'')
\]
是\Term{exact}{正合}的。
\end{SGAresult}
由于 $A'$ 在 $A$ 上\Term{faithfully-flat}{忠实平坦}，\Term{homomorphism}{同态} $N\to N'$ 为\Term{injectivity}{单射}，因此第一条箭头为\Term{injectivity}{单射}。只需证明：若 $A'$-\Term{homomorphism}{同态} $u':M'\to N'$ 与\Term{descent-datum}{下降数据}相容，则它来自一个 $A$-\Term{homomorphism}{同态} $u:M\to N$。这也就是要说，$u'$ 把 $M'$ 的子集 $M$ 映入 $N'$ 的子集 $N$；因为 $u'$ 为 $A'$-线性，所诱导的 $u:M\to N$ 自动为 $A$-线性，也同样可见必有 $u'=u\otimes_A A'$。而对 $x\in M$，$u'(x)$ 属于映射对 $N'\rightrightarrows N''$ 的核。因此，要证明 1.4，只需证明以下对应于 $M=A$ 的特例。

\begin{SGAresult}{推论}{1.5}
设 $N$ 为 $A$-模，则集合映射序列
\[
N\longrightarrow N'\rightrightarrows N''
\]
\Term{exact}{正合}。
\end{SGAresult}
事实上，设 $A_1$ 为\Term{faithfully-flat}{忠实平坦}的 $A$-代数。要证明上述序列\Term{exact}{正合}，只需证明沿 $A\to A_1$ 作环变换后所得的序列\Term{exact}{正合}。不难看出，后者正是由 $A_1$-模 $N_1=N\otimes_A A_1$ 及 $A_1$-代数 $A'_1=A_1\otimes_A A'$ 得到的序列。因此，只需找到一个在 $A$ 上\Term{faithfully-flat}{忠实平坦}的 $A_1$，使 $\operatorname{Spec}(A'_1)\to\operatorname{Spec}(A_1)$ 为严格的 $\mathcal F$-\Term{descent-morphism}{下降态射}。取 $A_1=A'$ 即可：此时该\Term{morphism}{态射}有右逆，故由 [D]，对 $\operatorname{Sch}$ 上任意\Term{fibred-category}{纤维范畴}，它都是一个\Term{effective-descent-morphism}{有效下降态射}。

最后，设 $N'$ 是带有相对于 $A\to A'$ 的\Term{descent-datum}{下降数据}的 $A'$-模，即赋有一个\Term{isomorphism}{同构}
\[
\varphi:N'_1\xrightarrow{\sim}N'_2,
\]
这里两端是由 $N$ 沿两个环变换 $A'\rightrightarrows A'\otimes_A A'$ 得到的模。还需证明：连同其\Term{descent-datum}{下降数据}，$N'$ \Term{isomorphism}{同构}于某个 $N\otimes_A A'$。利用 1.5，不难看出这一断言等价于以下引理。

\SourcePage{172}
\begin{SGAresult}{引理}{1.6}
设 $N'$ 为带有相对于 $A\to A'$ 的\Term{descent-datum}{下降数据}的 $A'$-模，其中 $A'$ 是一个 $A$-代数。令 $N$ 为 $N'$ 中所有满足
\[
\varphi(x\otimes_{A'}1_{A'})=1_{A'}\otimes_{A'}x
\]
的 $x$ 所成的 $A$-子模，并考虑典范\Term{homomorphism}{同态}
\[
N\otimes_A A'\longrightarrow N'.
\]
它与\Term{descent-datum}{下降数据}相容。若 $A'$ 在 $A$ 上\Term{faithfully-flat}{忠实平坦}，则这个\Term{homomorphism}{同态}为\Term{isomorphism}{同构}。
\end{SGAresult}
证明如下。仍取一个\Term{faithfully-flat}{忠实平坦}的 $A$-代数 $A_1$；要证明所考虑的\Term{morphism}{态射}是\Term{isomorphism}{同构}，只需证明在环变换 $A\to A_1$ 后它成为\Term{isomorphism}{同构}。利用 $A_1$ 在 $A$ 上的\Term{flat}{平坦}性，可以看出所得\Term{homomorphism}{同态}，正是直接从 $A'_1=A'\otimes_A A_1$ 上的模 $N'\otimes_A A_1$ 得到的那个\Term{homomorphism}{同态}；此模带有相对于 $A_1\to A'_1$ 的\Term{descent-datum}{下降数据}，是由 $N'$ 上原有的\Term{descent-datum}{下降数据}作环变换典范地得到的。因此，只需找到一个在 $A$ 上\Term{faithfully-flat}{忠实平坦}的 $A_1$，使 $\operatorname{Spec}(A'_1)\to\operatorname{Spec}(A_1)$ 为 $\mathcal F$-\Term{effective-descent-morphism}{有效下降态射}。如前取 $A_1=A'$ 即可。引理 1.6 得证，定理 1.1 的证明也随之完成。

\begin{SGAresult}{推论}{1.7（模\Term{section}{截面}的\Term{descent}{下降}）}
设 $g:S'\to S$ 是\Term{faithfully-flat}{忠实平坦}且\Term{quasi-compact}{拟紧}的\Term{prescheme}{预概形}\Term{morphism}{态射}。对 $S$ 上任意\Term{quasi-coherent-module}{拟凝聚模} $G$，记 $G',G''$ 为它在 $S'$ 和 $S''=S'\times_S S'$ 上的\Term{inverse-image}{逆像}，考虑 $S$ 上模\Term{homomorphism}{同态}的图
\[
G\longrightarrow g_*(G')\rightrightarrows h_*(G''),
\]
其中 $h:S''\to S$ 是结构\Term{morphism}{态射}。这个图是\Term{exact}{正合}的。
\end{SGAresult}
事实上，这意味着：对 $S$ 中每个开集 $U$，由 $U$ 上的\Term{section}{截面}组成的相应图\Term{exact}{正合}。显然可假设 $U=S$，而此时所需的\Term{exactness}{正合性}就是在 1.2 中令 $F=\mathcal O_S$ 得到的特例。

由于模的\Term{inverse-image}{逆像}\Term{functor}{函子}\Term{right-exact}{右正合}，从 1.1 形式地推出：
\begin{SGAresult}{推论}{1.8（商模的\Term{descent}{下降}）}
沿用 1.7 的记号。对任一\Term{prescheme}{预概形}上的\Term{quasi-coherent-module}{拟凝聚模} $F$，以 $\operatorname{Quot}(F)$ 表示 $F$ 的\Term{quasi-coherent}{拟凝聚}商模所成的集合。依此约定，集合映射图
\[
\operatorname{Quot}(G)\longrightarrow
\operatorname{Quot}(G')\rightrightarrows
\operatorname{Quot}(G'')
\]
\Term{exact}{正合}。
\end{SGAresult}
显然，把商模换成子模，也有同样的结论，因为两者一一对应。特别令 $G=\mathcal O_S$，就得到：

\SourcePage{173}
\begin{SGAresult}{推论}{1.9（闭子\Term{prescheme}{预概形}的\Term{descent}{下降}）}
对每个\Term{prescheme}{预概形} $X$，以 $H(X)$ 表示其闭子\Term{prescheme}{预概形}的集合。在 1.7 的条件下，集合映射图
\[
H(S)\longrightarrow H(S')\rightrightarrows H(S'')
\]
\Term{exact}{正合}。
\end{SGAresult}
还应以下列结果补充定理 1.1。

\begin{SGAresult}{命题}{1.10（模的性质的\Term{descent}{下降}）}
设 $g:S'\to S$ \Term{faithfully-flat}{忠实平坦}且\Term{quasi-compact}{拟紧}，$F$ 为 $S$ 上的\Term{quasi-coherent-module}{拟凝聚模}。则 $F$ 为\Term{finite-type}{有限型}、\Term{finite-presentation}{有限表示}、\Term{locally-free}{局部自由}且\Term{finite-type}{有限型}，当且仅当其在 $S'$ 上的\Term{inverse-image}{逆像} $F'$ 分别有相应性质。
\end{SGAresult}
只需证明充分性。显然可假设 $S$ \Term{affine}{仿射}；再把 $S'$ 换成一族\Term{covering}{覆盖}它的\Term{affine-open}{仿射开集}的不交并，便归结为 $S'$ 也\Term{affine}{仿射}的情形。此时结论等价于下面的断言。

\begin{SGAresult}{推论}{1.11}
设 $A$ 为环，$A'$ 为\Term{faithfully-flat}{忠实平坦}的 $A$-代数，$M$ 为 $A$-模，$M'=M\otimes_A A'$。则 $M$ 为\Term{finite-type}{有限型}（分别为\Term{finite-presentation}{有限表示}、\Term{locally-free}{局部自由}\Term{finite-type}{有限型}），当且仅当 $M'$ 有相应性质。
\end{SGAresult}
事实上，有 $M=\varinjlim_i M_i$，其中 $M_i$ 遍历 $M$ 的\Term{finite-type}{有限型}子模。于是 $M'=\varinjlim_i M'_i$；若 $M'$ 为\Term{finite-type}{有限型}，它就等于某个 $M'_i$，故由\Term{faithfully-flat}{忠实平坦}性，$M=M_i$，从而 $M$ 为\Term{finite-type}{有限型}。因此存在\Term{exact-sequence}{正合列}
\[
0\longrightarrow R\longrightarrow L\longrightarrow M\longrightarrow0,
\]
其中 $L$ 为\Term{finite-type}{有限型}\Term{free-module}{自由模}，由此得到\Term{exact-sequence}{正合列}
\[
0\longrightarrow R'\longrightarrow L'\longrightarrow M'\longrightarrow0,
\]
其中 $L'$ 为\Term{finite-type}{有限型}\Term{free-module}{自由模}。因此，若 $M'$ 为\Term{finite-presentation}{有限表示}，$R'$ 就是\Term{finite-type}{有限型}的；由前面的结论，$R$ 为\Term{finite-type}{有限型}，故 $M$ 为\Term{finite-presentation}{有限表示}。最后，$M$ \Term{locally-free}{局部自由}且\Term{finite-type}{有限型}，等价于它\Term{finite-presentation}{有限表示}且\Term{flat}{平坦}；\Term{noetherian}{诺特}情形见第 IV 讲，一般情形留给读者。由于这两个性质都能\Term{descent}{下降}，它们的合取也能\Term{descent}{下降}，证明完毕。

\textbf{注记 1.12。}由 1.1 与 1.10 合起来可知，将\Term{fibred-category}{纤维范畴} $\mathcal F$ 换成由\Term{finite-type}{有限型}、\Term{finite-presentation}{有限表示}、\Term{locally-free}{局部自由}\Term{finite-type}{有限型}或给定秩 $n$ 的\Term{locally-free}{局部自由}\Term{quasi-coherent-module}{拟凝聚模}组成的\Term{fibre}{纤维}子\Term{category}{范畴}，1.1 的结论仍然成立。

\SourcePage{174}
\SGAsection{2}{在另一个预概形上仿射的预概形的下降}
由于模的\Term{inverse-image}{逆像}\Term{functor}{函子}与\Term{tensor-product}{张量积}及其他张量运算相容，定理 1.1 有多种变式：把单个\Term{quasi-coherent-module}{拟凝聚模}换成一个\Term{quasi-coherent-module}{拟凝聚模}，或一族\Term{quasi-coherent-module}{拟凝聚模}，并赋予可以用张量运算表述的附加结构即可。例如，给定 $S$ 上三个\Term{quasi-coherent-module}{拟凝聚模} $F,G,H$ 及一个\Term{pairing}{配对}
\[
F\otimes G\longrightarrow H,
\]
等价于给定 $S'$ 上三个\Term{quasi-coherent-module}{拟凝聚模} $F',G',H'$，赋予相对于 $g:S'\to S$ 的\Term{descent-datum}{下降数据}，并给定一个按显然意义与这些\Term{descent-datum}{下降数据}“相容”的\Term{pairing}{配对}
\[
F'\otimes G'\longrightarrow H'.
\]
例如，令 $F=G=H$，便可看出：给定 $S$ 上的\Term{quasi-coherent-module}{拟凝聚模} $F$ 及一个代数乘法——暂时不要求结合律、交换律或单位\Term{section}{截面}的存在——等价于在 $S'$ 上给定同样的对象，再赋予\Term{descent-datum}{下降数据}。利用上一节的结果，立刻可见，$F$ 满足刚才提到的任一通常公理，当且仅当 $F'$ 满足相应公理。例如，给定 $S$ 上一个\Term{quasi-coherent}{拟凝聚}代数 $\mathcal A$——从现在起，这一说法默认它结合、交换且有单位\Term{section}{截面}——等价于给定 $S'$ 上一个\Term{quasi-coherent}{拟凝聚}代数 $\mathcal A'$，并赋予相对于 $g:S'\to S$ 的\Term{descent-datum}{下降数据}。回顾 $S$ 上\Term{quasi-coherent}{拟凝聚}代数的\Term{opposite-category}{反范畴}与 $S$ 上\Term{affine}{仿射}的 $S$-\Term{prescheme}{预概形}\Term{category}{范畴}之间的等价（EGA II 第 1 节），立即得到：

\begin{SGAresult}{定理}{2.1}
设 $\mathcal F'$ 为\Term{prescheme}{预概形}\Term{affine-morphism}{仿射态射} $f:X\to S$ 的\Term{fibred-category}{纤维范畴}，把它看成\Term{prescheme}{预概形}\Term{category}{范畴} $\operatorname{Sch}$ 中箭头的\Term{fibred-category}{纤维范畴}的\Term{fibre}{纤维}子\Term{category}{范畴}（第 VI 讲 11.a）。设 $g:S'\to S$ 为\Term{faithfully-flat}{忠实平坦}且\Term{quasi-compact}{拟紧}的\Term{prescheme}{预概形}\Term{morphism}{态射}。则 $g$ 是 $\mathcal F'$-\Term{effective-descent-morphism}{有效下降态射}。
\end{SGAresult}

\SGAsection{3}{态射的集合性质与有限性性质的下降}
\footnote{关于第 3、4 节所涉及的其他结果，参见 EGA IV 2.3、2.6、2.7。}
\begin{SGAresult}{命题}{3.1}
设 $f:X\to Y$ 为 $S$-\Term{morphism}{态射}，$g:S'\to S$ 为\Term{surjectivity}{满射}，$f':X'=X\times_S S'\to Y'=Y\times_S S'$ 为沿 $g:S'\to S$ 作\Term{base-change}{基变换}得到的\Term{morphism}{态射}。则 $f$ 为\Term{surjectivity}{满射}（分别为\Term{radicial}{根性}\Term{morphism}{态射}），当且仅当 $f'$ 有相应性质。
\end{SGAresult}

\SourcePage{175}
注意，$f'$ 也可由\Term{base-change}{基变换} $Y'\to Y$ 得到，而后者同样为\Term{surjectivity}{满射}，因为它由\Term{surjectivity}{满射} $g:S'\to S$ 作\Term{base-change}{基变换}而来。另一方面，对每个 $y\in Y$ 及位于 $y$ 上方的 $y'\in Y'$，有\Term{isomorphism}{同构}
\[
X'_{y'}\xleftarrow{\sim}X_y\otimes_{k(y)}k(y'),
\]
其中 $X_y$ 表示 $X$ 在 $y$ 处的\Term{fibre}{纤维}，$X'_{y'}$ 表示 $X'$ 在 $y'$ 处的\Term{fibre}{纤维}。因此，$X_y$ 非空（分别至多只有一个点，且该点所对应的\Term{residue-field}{剩余域}\Term{group-extension}{扩张}为\Term{purely-inseparable-extension}{纯不可分扩张}），当且仅当 $X'_{y'}$ 有同样性质。这就证明了 3.1。

\begin{SGAresult}{推论}{3.2}
在 3.1 的条件下，若 $f'$ 为\Term{injectivity}{单射}（分别为\Term{bijection}{双射}），则 $f$ 也如此。
\end{SGAresult}
这是因为，若 $X'_{y'}$ 至多有一个点（分别恰有一个点），则 $X_y$ 也如此；确实，\Term{morphism}{态射} $X'_{y'}\to X_y$ 是\Term{surjectivity}{满射}，因为它由\Term{surjectivity}{满射} $\operatorname{Spec}(k(y'))\to\operatorname{Spec}(k(y))$ 作\Term{base-change}{基变换}得到。

\begin{SGAresult}{命题}{3.3}
沿用 3.1 的记号。假设 $g:S'\to S$ 为\Term{surjectivity}{满射}且\Term{quasi-compact}{拟紧}（分别为\Term{faithfully-flat}{忠实平坦}且\Term{quasi-compact}{拟紧}）。则 $f$ \Term{quasi-compact}{拟紧}（分别为\Term{finite-type}{有限型}），当且仅当 $f'$ 有相应性质。
\end{SGAresult}
只需证明充分性。显然可假设 $S=Y$，因为对 $g:S'\to S$ 的假设对 $Y'\to Y$ 仍成立。还可假设 $Y$ \Term{affine}{仿射}。于是 $Y'$ \Term{quasi-compact}{拟紧}，故 $X'$ 也\Term{quasi-compact}{拟紧}，因为按假设 $f'$ \Term{quasi-compact}{拟紧}。设 $(X_i)_{i\in I}$ 为 $X$ 的一个\Term{affine}{仿射}开\Term{covering}{覆盖}；相应的 $X'_i$ 是\Term{covering}{覆盖} $X'$ 的开集，因此其中有一个\Term{finite}{有限}子族\Term{covering}{覆盖} $X'$。由于 $X'\to X$ 为\Term{surjectivity}{满射}，相应的 $X_i$ 已\Term{covering}{覆盖} $X$，从而 $X$ \Term{quasi-compact}{拟紧}，即 $f$ \Term{quasi-compact}{拟紧}。现在设 $f'$ 为\Term{finite-type}{有限型}，并假设 $g$ \Term{faithfully-flat}{忠实平坦}，来证明 $f$ 为\Term{finite-type}{有限型}。把 $Y'$ 换成一族\Term{covering}{覆盖}它的\Term{affine-open}{仿射开集}的不交并，可假设 $Y'$ \Term{affine}{仿射}。最后，由上述结论，$X$ 被\Term{finite}{有限}多个\Term{affine-open}{仿射开集} $X_i$ \Term{covering}{覆盖}；已知 $X'_i$ 在 $Y'$ 上\Term{finite-type}{有限型}，要证明 $X_i$ 在 $Y$ 上\Term{finite-type}{有限型}。这就归结为：

\begin{SGAresult}{推论}{3.4}
设 $B$ 为 $A$-代数，$A'$ 为\Term{faithfully-flat}{忠实平坦}的 $A$-代数，$B'=B\otimes_A A'$ 为 $B$ 作环变换所得的 $A'$-代数。则 $B$ 为\Term{finite-type}{有限型}，当且仅当 $B'$ 为\Term{finite-type}{有限型}。
\end{SGAresult}
只需证明充分性。有 $B=\varinjlim_i B_i$，其中 $B_i$ 为 $B$ 的\Term{finite-type}{有限型}子代数。因此 $B'=\varinjlim_i B'_i$；若 $B'$ 在 $A'$ 上\Term{finite-type}{有限型}，则它等于某个 $B'_i$，故 $B=B_i$，所以 $B$ 为\Term{finite-type}{有限型}。

\begin{SGAresult}{推论}{3.5}
仍假设\Term{base-change}{基变换}\Term{morphism}{态射} $g:S'\to S$ \Term{faithfully-flat}{忠实平坦}且\Term{quasi-compact}{拟紧}。则 $f$ 为\Term{quasi-finite}{拟有限}，当且仅当 $f'$ 为\Term{quasi-finite}{拟有限}。
\end{SGAresult}

\SourcePage{176}
事实上，按定义，“\Term{quasi-finite}{拟有限}”就是“\Term{finite-type}{有限型}”与“\Term{fibre}{纤维}\Term{finite}{有限}”两项性质的合取；两者都能沿 $g$ \Term{descent}{下降}，前者由 3.3，后者由 3.1 的论证，其中只使用了 $g$ 为\Term{surjectivity}{满射}这一事实。

\textbf{注记 3.6。}设 $A$ 为环，$X$ 为 $A$-\Term{prescheme}{预概形}。不难看出，以下条件等价：
\begin{enumerate}[label=(\roman*)]
\item 存在一个\Term{noetherian-ring}{诺特环} $A_0$——如有需要，可以取为 $A$ 的一个\Term{finite-type}{有限型}子环——、一个\Term{finite-type}{有限型} $A_0$-\Term{prescheme}{预概形} $X_0$、一个\Term{homomorphism}{同态} $A_0\to A$，以及一个 $A$-\Term{isomorphism}{同构} $X\xrightarrow{\sim}X_0\times_{A_0}A$。
\item \Term{diagonal-morphism}{对角态射} $X\to X\times_{\operatorname{Spec}(A)}X$ \Term{quasi-compact}{拟紧}；若 $X$ 在 $A$ 上\Term{separated}{分离}，这一条件自动成立。另外，$X$ 为\Term{finite}{有限}多个\Term{affine-open}{仿射开集} $X_i$ 的并，其环 $B_i$ 是 $A$ 上的\Term{finite-presentation}{有限表示}代数，即\Term{finite}{有限}个不定元的多项式代数模去一个\Term{finite-type}{有限型}理想所得的商。
\end{enumerate}
若 $X$ 本身\Term{affine}{仿射}，环为 $B$，则这些条件只意味着 $B$ 是 $A$ 上的\Term{finite-presentation}{有限表示}代数。

若 $Y$ 可由\Term{affine-open}{仿射开集} $Y_i$ \Term{covering}{覆盖}，使 $X|_{Y_i}$ 作为 $Y_i$-\Term{prescheme}{预概形}满足上述等价条件，就称 $f:X\to Y$ 为\Term{finite-presentation}{有限表示}\Term{morphism}{态射}，也说 $X$ 在 $Y$ 上\Term{finite-presentation}{有限表示}。此时对 $Y$ 中任意\Term{affine-open}{仿射开集} $Y'$，$X|_{Y'}$ 也满足这些条件。这一性质在\Term{base-change}{基变换}下保持，而且两个\Term{finite-presentation}{有限表示}\Term{morphism}{态射}的复合仍为\Term{finite-presentation}{有限表示}。

有了这些概念后，由 (ii) 按 1.10 的方式论证，可以看出：把该断言中的“\Term{finite-type}{有限型}”换成“\Term{finite-presentation}{有限表示}”，结论仍成立。

\SGAsection{4}{拓扑性质的下降}
\begin{SGAresult}{定理}{4.1}
设 $g:Y'\to Y$ 为\Term{morphism}{态射}，$Z$ 为 $Y$ 的子集。假设 $g$ \Term{flat}{平坦}，且存在\Term{quasi-compact}{拟紧}\Term{morphism}{态射} $f:X\to Y$ 使 $Z=f(X)$。注意，若 $Y$ 为\Term{noetherian}{诺特}的，则“$Z$ \Term{constructible}{可构造}”蕴含后一条件。于是
\[
g^{-1}(\overline Z)=\overline{g^{-1}(Z)}.
\]
\end{SGAresult}
可先假设 $Y$ \Term{affine}{仿射}，再假设 $Y'$ \Term{affine}{仿射}。由于 $Y$ \Term{affine}{仿射}，$X$ 为\Term{finite}{有限}多个\Term{affine-open}{仿射开集} $X_i$ 的并；把 $X$ 换成这些 $X_i$ 的不交并，可进一步假设 $X$ \Term{affine}{仿射}。设 $A,A',B$ 分别为 $Y,Y',X$ 的环，$B'=B\otimes_A A'$ 为 $X'=X\times_Y Y'$ 的环，$I$ 为 $A\to B$ 的核，$I'$ 为 $A'\to B'$ 的核。它们在 $Y,Y'$ 中定义的闭集，分别是 $Z=f(X)$ 的闭包及 $Z'=f'(X')=g^{-1}(Z)$ 的闭包。要证明后者等于 $g^{-1}(\overline Z)$，只需证明 $I'=IA'$，而这由 $A'$ 在 $A$ 上的\Term{flat}{平坦}性立即得到。

\SourcePage{177}
\begin{SGAresult}{推论}{4.2}
设 $g:Y'\to Y$ 为\Term{flat}{平坦}且\Term{quasi-compact}{拟紧}的\Term{morphism}{态射}，$Z'$ 为 $Y'$ 中对 $g$ 所定义的集合意义下的等价关系饱和的闭子集。则
\[
Z'=g^{-1}\bigl(\overline{g(Z')}\bigr).
\]
\end{SGAresult}
事实上，令 $Z=g(Z')$，则 $Z'=g^{-1}(Z)$。可以应用 4.1：取 $X$ 为带有 $Y'$ 所诱导的既约结构的\Term{prescheme}{预概形} $Z'$，就满足该定理对 $Z$ 的条件。这里 $g$ \Term{quasi-compact}{拟紧}，保证了诱导\Term{morphism}{态射} $f:Z'\to Y$ \Term{quasi-compact}{拟紧}。

断言 4.2 也意味着，$g(Y')$ 由 $Y$ 诱导的拓扑，是 $Y'$ 的拓扑的\Term{quotient-topology}{商拓扑}。特别地：
\begin{SGAresult}{推论}{4.3}
设 $g:Y'\to Y$ \Term{faithfully-flat}{忠实平坦}且\Term{quasi-compact}{拟紧}。则 $g$ 使 $Y$ 成为 $Y'$ 的\Term{quotient-topology}{商拓扑}空间。换言之，对 $Y$ 的子集 $Z$，$Z$ 为闭集（分别为开集），当且仅当 $Z'=g^{-1}(Z)$ 有相应性质。
\end{SGAresult}
现在回顾，$Y'$ 的两个元素 $a,b$ 在 $Y$ 中有相同的像，当且仅当存在合适的 $c\in Y''=Y'\times_Y Y'$，使 $a=p_1(c),b=p_2(c)$。因此，若 $g$ 为\Term{surjectivity}{满射}，就有集合的\Term{exact}{正合}图
\[
\mathcal P(Y)\longrightarrow\mathcal P(Y')\rightrightarrows\mathcal P(Y''),
\]
其中对任意集合 $E$，$\mathcal P(E)$ 表示其幂集。于是 4.3 也\Term{representable}{可表}述为：
\begin{SGAresult}{推论}{4.4（开子集及闭子集的\Term{descent}{下降}）}
设 $g:Y'\to Y$ 如 4.3。对任意\Term{prescheme}{预概形} $X$，分别以 $\operatorname{Ouv}(X)$ 和 $\operatorname{Fer}(X)$ 表示其开子集和闭子集的集合。由 $g$ 及 $Y''=Y'\times_Y Y'$ 的两个投影，得到以下集合映射\Term{exact}{正合}图：
\[
\operatorname{Ouv}(Y)\longrightarrow\operatorname{Ouv}(Y')\rightrightarrows\operatorname{Ouv}(Y''),
\]
\[
\operatorname{Fer}(Y)\longrightarrow\operatorname{Fer}(Y')\rightrightarrows\operatorname{Fer}(Y'').
\]
\end{SGAresult}
还有对 4.3 的如下补充。
\begin{SGAresult}{推论}{4.5}
设 $g:Y'\to Y$ 如 4.3，$Z$ 为 $Y$ 的子集，且存在像为 $Z$ 的\Term{quasi-compact}{拟紧}\Term{morphism}{态射} $f:X\to Y$；例如，$Y$ 为\Term{noetherian}{诺特}的且 $Z$ \Term{constructible}{可构造}。则 $Z$ 是 $Y$ 的\Term{locally-closed}{局部闭}子集，当且仅当 $Z'=g^{-1}(Z)$ 是 $Y'$ 的\Term{locally-closed}{局部闭}子集。
\end{SGAresult}
只需证明充分性。设 $Y_1$ 为 $Y$ 的闭子\Term{prescheme}{预概形}，其底集为 $Z$ 的闭包，结构为诱导的既约结构；设 $Y'_1=Y_1\times_Y Y'$ 为其\Term{inverse-image}{逆像}闭子\Term{prescheme}{预概形}。其底集是 $g^{-1}(Y_1)=g^{-1}(\overline Z)$，由 4.1，这等于 $\overline{Z'}$。由于 $Z'$ 在 $Y'$ 中\Term{locally-closed}{局部闭}，它在 $\overline{Z'}$ 中开，故在 $Y'_1$ 中开。后者在 $Y_1$ 上\Term{faithfully-flat}{忠实平坦}且\Term{quasi-compact}{拟紧}，故由 4.3，$Z$ 在 $Y_1$ 中开，即在 $\overline Z$ 中开；这就说明 $Z$ \Term{locally-closed}{局部闭}。

\SourcePage{178}
\begin{SGAresult}{推论}{4.6}
设 $g:S'\to S$ 为\Term{faithfully-flat}{忠实平坦}且\Term{quasi-compact}{拟紧}的\Term{morphism}{态射}，$f:X\to Y$ 为 $S$-\Term{morphism}{态射}，$f':X'\to Y'$ 为由它作\Term{base-change}{基变换}得到的 $S'$-\Term{morphism}{态射}。假设 $f'$ 是开映射（分别为闭映射、\Term{quasi-compact}{拟紧}且为到其像的同胚、满同胚），则 $f$ 有相应性质。
\end{SGAresult}
由于 $Y'$ 在 $Y$ 上\Term{faithfully-flat}{忠实平坦}且\Term{quasi-compact}{拟紧}，可以假设 $Y=S$。设 $Q$ 为 $X$ 的子集，以 $h$ 表示投影\Term{morphism}{态射} $X'\to X$，则
\[
g^{-1}(f(Q))=f'(h^{-1}(Q)).
\]
若 $Q$ 开（分别闭），则 $h^{-1}(Q)$ 也开（分别闭）。若 $f'$ 为开映射（分别闭映射），则 $f'(h^{-1}(Q))$ 也如此；由上式及 4.3，$f(Q)$ 便有相同性质。这就证明了前两个断言。还需讨论 $f'$ 为到其像的同胚的情形，并证明 $f$ 也是到其像的同胚；满同胚的情形随后由 3.1 得到。由 3.2，$f$ 为\Term{injectivity}{单射}，只需证明映射 $X\to f(X)$ 为开映射。又由 3.3，已知 $f$ \Term{quasi-compact}{拟紧}。只需证明，对 $X$ 的任意闭子集 $Z$，有
\[
Z=f^{-1}\bigl(\overline{f(Z)}\bigr).
\]
这等价于沿\Term{surjectivity}{满射} $h:X'\to X$ 取\Term{inverse-image}{逆像}后所得的相应公式，即令 $Z'=h^{-1}(Z)$，则
\[
Z'=f'^{-1}\bigl(g^{-1}(\overline{f(Z)})\bigr).
\]
由 4.1 对 $Y$ 的子集 $f(Z)$ 应用所得的公式
\[
g^{-1}(\overline{f(Z)})=\overline{g^{-1}(f(Z))},
\]
待证公式等价于
\[
Z'=f'^{-1}\bigl(\overline{f'(Z')}\bigr),
\]
这由 $f'$ 为到其像的同胚的假设成立。

注意：最后的论证在已假设 $f$ \Term{quasi-compact}{拟紧}之后，没有用到 $g$ \Term{quasi-compact}{拟紧}，只用了 $g$ \Term{faithfully-flat}{忠实平坦}。因此，在这一假设下，就能\Term{descent}{下降}“到其像的同胚”或“满同胚”的性质；由前面的论证，还能\Term{descent}{下降}“$f'$ \Term{quasi-compact}{拟紧}，且使 $f'(X')$ 成为 $X'$ 的\Term{quotient-topology}{商拓扑}空间”这一性质。

若\Term{prescheme}{预概形}\Term{morphism}{态射} $f:X\to Y$ 对任意\Term{base-change}{基变换} $Y'\to Y$，所得 $f':X'\to Y'$ 都是开映射（分别为闭映射、到其像空间的同胚），就称 $f$ \Term{universally-open}{普遍开}（分别为\Term{universally-closed}{普遍闭}、普遍双连续等）。由 4.6 得到：
\begin{SGAresult}{推论}{4.7}
在 4.6 的条件下，$f$ \Term{universally-open}{普遍开}（分别为\Term{universally-closed}{普遍闭}、普遍为到其像的同胚、\Term{universal-homeomorphism}{普遍同胚}），当且仅当 $f'$ 有相应性质。
\end{SGAresult}

\SourcePage{179}
\begin{SGAresult}{推论}{4.8}
在 4.6 的条件下，$f$ 为\Term{separated-morphism}{分离态射}（分别为\Term{proper-morphism}{固有态射}），当且仅当 $f'$ 有相应性质。
\end{SGAresult}
$f$ \Term{separated}{分离}，就是说\Term{diagonal-morphism}{对角态射} $X\to X\times_Y X$ 为闭映射，也可说为\Term{universally-closed}{普遍闭}映射，因此第一个断言由 4.7 得到。$f$ \Term{proper}{固有}，就是说它满足：a) $f$ \Term{finite-type}{有限型}；b) $f$ \Term{separated}{分离}；c) $f$ \Term{universally-closed}{普遍闭}。性质 a) 由 3.3 可以\Term{descent}{下降}，b) 由刚证明的结论可以\Term{descent}{下降}，c) 也由 4.7 可以\Term{descent}{下降}。

\textbf{注记 4.9。}回顾一下，若 $g:Y'\to Y$ 为\Term{finite-type}{有限型}\Term{flat}{平坦}\Term{morphism}{态射}，且 $Y$ \Term{locally-noetherian}{局部诺特}，则 $g$ 为\Term{open-morphism}{开态射}（VI 6.6）；这是比 4.3 更精确的结果。不过应当注意，\Term{noetherian}{诺特}\Term{prescheme}{预概形}之间\Term{faithfully-flat}{忠实平坦}且\Term{quasi-compact}{拟紧}的\Term{morphism}{态射} $f$ 一般并不是\Term{open-morphism}{开态射}。例如，设 $Y$ 为\Term{irreducible}{不可约}\Term{scheme}{概形}，其一般点 $y$ 不是开点，如一条代数曲线；取 $Y'=Y\amalg\operatorname{Spec}(k(y))$，则开子集 $\operatorname{Spec}(k(y))$ 在结构\Term{morphism}{态射} $Y'\to Y$ 下的像，不是 $Y$ 的开子集。读者还会注意到，若取消所考虑的\Term{faithfully-flat}{忠实平坦}\Term{morphism}{态射}同时\Term{quasi-compact}{拟紧}的假设，本讲的若干断言就不再成立。典型反例是：取 $Y'$ 为 $Y$ 的所有点的\Term{local-ring}{局部环}的谱的不交并。例如，仍取 $Y$ 为\Term{irreducible}{不可约}代数曲线，令 $Z$ 为仅含其一般点的子集，则 $Z$ 在 $Y'$ 中的\Term{inverse-image}{逆像}是开的，但 $Z$ 本身并不开。

\textbf{4.10。}本讲若干断言中，可将“$Y'$ 在 $Y$ 上\Term{flat}{平坦}”替换为：存在 $Y'$ 上的\Term{finite-type}{有限型}模 $F$，其\Term{support}{支集}为 $Y'$，且 $F$ 相对于 $Y$ \Term{flat}{平坦}。相应地，把\Term{faithfully-flat}{忠实平坦}性的假设换成上述条件，再加上 $Y'\to Y$ 为\Term{surjectivity}{满射}。这适用于 1.10 的前两个断言、3.3、3.5、4.1，从而也适用于本节的全部结果。

\SGAsection{5}{预概形态射的下降}
\begin{SGAresult}{命题}{5.1}
设 $g:S'\to S$ 为\Term{prescheme}{预概形}\Term{morphism}{态射}。

\noindent a) 假设 $g$ 为\Term{surjectivity}{满射}，且\Term{homomorphism}{同态}
\[
g^*:\mathcal O_S\longrightarrow g_*(\mathcal O_{S'})
\]
为\Term{injectivity}{单射}，则 $g$ 是\Term{prescheme}{预概形}\Term{category}{范畴}中的\Term{epimorphism}{满态射}，甚至是环空间\Term{category}{范畴}中的\Term{epimorphism}{满态射}。
\SourcePage{180}

\noindent b) 假设 $g$ 为\Term{surjectivity}{满射}，并使 $S$ 成为 $S'$ 的\Term{quotient-topology}{商拓扑}空间。令 $S''=S'\times_S S'$，$h:S''\to S$ 为结构\Term{morphism}{态射}，考虑典范\Term{homomorphism}{同态}图
\[
\mathcal O_S\longrightarrow g_*(\mathcal O_{S'})\rightrightarrows h_*(\mathcal O_{S''}),
\]
并假设它\Term{exact}{正合}。则 $g$ 是\Term{prescheme}{预概形}\Term{category}{范畴}中的\Term{effective-epimorphism}{有效满态射}，也同样是环空间\Term{category}{范畴}中的\Term{effective-epimorphism}{有效满态射}，即图
\[
S\longleftarrow S'\leftleftarrows S''
\]
\Term{exact}{正合}。
\end{SGAresult}
\begin{proof}
a) 要证明：环空间\Term{morphism}{态射} $f:S\to Z$ 由 $fg$ 确定。由于 $g$ 为\Term{surjectivity}{满射}，$f$ 的底层集合映射 $f_0$ 已经确定，剩下的是确定 $f$ 所定义的环层\Term{homomorphism}{同态} $\mathcal O_Z\to\mathcal O_S$，等价地，确定环层\Term{homomorphism}{同态}
\[
u:f_0^{-1}(\mathcal O_Z)\longrightarrow\mathcal O_S.
\]
已知 $fg$ 所定义的\Term{homomorphism}{同态}
\[
(fg)_0^{-1}(\mathcal O_Z)=g_0^{-1}\bigl(f_0^{-1}(\mathcal O_Z)\bigr)\longrightarrow\mathcal O_{S'},
\]
等价地，已知一个\Term{homomorphism}{同态}
\[
f_0^{-1}(\mathcal O_Z)\longrightarrow g_{0*}(\mathcal O_{S'})=g_*(\mathcal O_{S'}).
\]
立即可见，后者正是 $u$ 与 $g^*:\mathcal O_S\to g_*(\mathcal O_{S'})$ 的复合。由于 $g^*$ 为\Term{injectivity}{单射}，$u$ 便由 $g^*u$ 确定。

〔注意：这里显然没有用到 $g:S'\to S$ 是\Term{prescheme}{预概形}\Term{morphism}{态射}；对于任意环空间\Term{morphism}{态射}，断言都成立。对 b) 也有同样的说明，无论是在环空间\Term{category}{范畴}中，还是在\Term{local-ring}{局部环}空间\Term{category}{范畴}中。此外，若\Term{prescheme}{预概形}\Term{morphism}{态射} $g$ 未必\Term{surjectivity}{满射}，但 $g^*:\mathcal O_S\to g_*(\mathcal O_{S'})$ 为\Term{injectivity}{单射}，则对 $S$ 到某个\Term{scheme}{概形} $Z$ 的两个\Term{morphism}{态射} $f_1,f_2$，由 $f_1g=f_2g$ 仍有 $f_1=f_2$。事实上，设 $I$ 为 $S$ 上定义 $f_1,f_2$ 的重合子\Term{prescheme}{预概形}的理想，即沿 $(f_1,f_2)$ 取 $Z\times Z$ 中对角子\Term{prescheme}{预概形}的\Term{inverse-image}{逆像}所得的理想，便有 $I\subset\operatorname{Ker}(g^*)$。〕

b) 要证明：对任意环空间 $Z$，集合映射图
\[
\operatorname{Hom}(S,Z)\longrightarrow\operatorname{Hom}(S',Z)\rightrightarrows\operatorname{Hom}(S'',Z)
\]
\Term{exact}{正合}；而当 $Z$ 为\Term{local-ring}{局部环}空间，且只考虑\Term{local-ring}{局部环}空间\Term{morphism}{态射}时，结论也成立。由 a) 已知第一条映射为\Term{injectivity}{单射}，因此只需证明：若环空间\Term{homomorphism}{同态} $f':S'\to Z$ 满足 $f'p_1=f'p_2$，则存在环空间\Term{homomorphism}{同态} $f:S\to Z$ 使 $f'=fg$。
\SourcePage{181}
由于 $g$ 为\Term{surjectivity}{满射}，显然若 $f'$ 是\Term{local-ring}{局部环}空间\Term{morphism}{态射}，则 $f$ 也如此。

由对 $f'$ 的假设，底层集合映射 $f'_0$ 在 $g_0$ 的各条\Term{fibre}{纤维}上为常值。由于 $g_0$ \Term{surjectivity}{满射}，$f'_0$ 唯一分解为 $f'_0=f_0g_0$；其中 $f_0:S\to Z$ 必连续，因为 $g_0$ 把 $S$ 等同于 $S'$ 的\Term{quotient-topology}{商拓扑}空间。再考虑由对应于 $f'$ 的\Term{homomorphism}{同态} $(f_0g_0)^{-1}(\mathcal O_Z)\to\mathcal O_{S'}$ 得到的\Term{homomorphism}{同态}
\[
f_0^{-1}(\mathcal O_Z)\longrightarrow g_*(\mathcal O_{S'}).
\]
假设 $f'p_1=f'p_2$，意味着将此\Term{homomorphism}{同态}分别与两个\Term{homomorphism}{同态} $g_*(\mathcal O_{S'})\rightrightarrows h_*(\mathcal O_{S''})$ 复合，所得\Term{homomorphism}{同态}相同。因此由 b) 的假设，它可经由一个\Term{morphism}{态射}
\[
f_0^{-1}(\mathcal O_Z)\longrightarrow\mathcal O_S
\]
分解。后者定义一个环空间\Term{morphism}{态射} $f:S\to Z$，即为所求。
\end{proof}

\begin{SGAresult}{定理}{5.2}
设 $\mathcal F$ 为\Term{prescheme}{预概形}\Term{category}{范畴} $\operatorname{Sch}$ 中箭头的\Term{fibred-category}{纤维范畴}（VI 11.a）。则每个\Term{faithfully-flat}{忠实平坦}且\Term{quasi-compact}{拟紧}的\Term{morphism}{态射} $g:S'\to S$ 都是 $\mathcal F$-\Term{descent-morphism}{下降态射}，也称为 $\operatorname{Sch}$ 中的\Term{descent-morphism}{下降态射}。
\end{SGAresult}
具体说来，令 $S''=S'\times_S S'$，对 $S$ 上的两个\Term{prescheme}{预概形} $X,Y$，分别考虑它们在 $S'$ 上的\Term{inverse-image}{逆像} $X',Y'$，以及在 $S''$ 上的\Term{inverse-image}{逆像} $X'',Y''$，得到集合映射图
\[
\operatorname{Hom}_S(X,Y)\longrightarrow
\operatorname{Hom}_{S'}(X',Y')\rightrightarrows
\operatorname{Hom}_{S''}(X'',Y'').
\]
定理 5.2 的意思就是此图\Term{exact}{正合}。注意，一般不能断言 $g$ 为\Term{effective-descent-morphism}{有效下降态射}，即不能说：对 $S'$ 上任意\Term{prescheme}{预概形} $X'$，其上相对于 $g:S'\to S$ 的每个\Term{descent-datum}{下降数据}都有效。有效性问题往往很微妙，将在第 7 节讨论。

由 [D] 的一般理论，并考虑到 $\operatorname{Sch}$ 中存在\Term{fibre}{纤维}积，已知断言 5.2 等价于：
\begin{SGAresult}{推论}{5.3}
\Term{faithfully-flat}{忠实平坦}且\Term{quasi-compact}{拟紧}的\Term{prescheme}{预概形}\Term{morphism}{态射}是\Term{universal-effective-epimorphism}{普遍有效满态射}。
\end{SGAresult}
由于\Term{faithfully-flat}{忠实平坦}且\Term{quasi-compact}{拟紧}的\Term{morphism}{态射}在任意\Term{base-extension}{基扩张}后仍有这两个性质，只需证明它为\Term{effective-epimorphism}{有效满态射}。应用准则 5.1 b)，再结合 4.3 与 1.7，就得到所需结论。

\begin{SGAresult}{推论}{5.4}
设 $g:S'\to S$ \Term{faithfully-flat}{忠实平坦}且\Term{quasi-compact}{拟紧}，$f:X\to Y$ 为 $S$-\Term{morphism}{态射}，$f':X'\to Y'$ 为沿\Term{base-change}{基变换} $S'\to S$ 得到的 $S'$-\Term{morphism}{态射}。
\SourcePage{182}
则 $f$ 是\Term{isomorphism}{同构}，当且仅当 $f'$ 是\Term{isomorphism}{同构}。
\end{SGAresult}
事实上，若 $f'$ 是\Term{isomorphism}{同构}，它对 $X',Y'$ 上的自然\Term{descent}{下降}结构也是\Term{isomorphism}{同构}。由 5.2，\Term{functor}{函子} $X\mapsto X'$ 从 $\operatorname{Sch}_{/S}$ 到带有相对于 $g$ 的\Term{descent-datum}{下降数据}的 $S'$-\Term{prescheme}{预概形}\Term{category}{范畴}是\Term{fully-faithful}{全忠实}的，因而得到结论。

\begin{SGAresult}{推论}{5.5}
在 5.4 的条件下，$f$ 是\Term{closed-immersion}{闭浸入}（分别为\Term{open-immersion}{开浸入}、\Term{quasi-compact}{拟紧}\Term{immersion}{浸入}），当且仅当 $f'$ 有相应性质。
\end{SGAresult}
如常可设 $Y=S$，且只需证明充分性。注意，$X'/Y'$ 带有相对于 $g:Y'\to Y$ 的\Term{descent-datum}{下降数据}，且结构\Term{morphism}{态射} $f':X'\to Y'$ 为\Term{immersion}{浸入}，从而是\Term{monomorphism}{单态射}；因此，由 $S''$ 到 $S'$ 的两个投影所得到的 $X'/Y'$ 在 $Y''$ 中的两个\Term{inverse-image}{逆像}\Term{subobject}{子对象}相同。若 $f'$ 为\Term{closed-immersion}{闭浸入}，由 1.9，存在 $Y$ 的闭子\Term{prescheme}{预概形} $X_1$，它沿 $g:Y'\to Y$ 的\Term{inverse-image}{逆像}为 $X'$。由相对于 $\mathcal F$-\Term{descent-morphism}{下降态射}的\Term{descent}{下降}问题之解的唯一性，$X_1$ 与 $X$ 为 $Y$-\Term{isomorphism}{同构}，故 $f:X\to Y$ 是\Term{closed-immersion}{闭浸入}。对\Term{open-immersion}{开浸入}利用 4.4 同样论证。最后，若 $f'$ 为\Term{quasi-compact}{拟紧}\Term{immersion}{浸入}，由 3.3，$f$ \Term{quasi-compact}{拟紧}，故可对 $Y$ 的子集 $f(X)$ 应用准则 4.5：由于它在 $Y'$ 中的\Term{inverse-image}{逆像} $f'(X')$ \Term{locally-closed}{局部闭}，$f(X)$ 也\Term{locally-closed}{局部闭}。再把 $Y$ 换成一个使 $f(X)$ 在其中闭的开子集，就归结为 $f'$ 是\Term{closed-immersion}{闭浸入}的情形，故由上述结论，$f$ 也是\Term{closed-immersion}{闭浸入}。

\begin{SGAresult}{推论}{5.6}
在 5.4 的条件下，$f$ 为\Term{affine-morphism}{仿射态射}，当且仅当 $f'$ 为\Term{affine-morphism}{仿射态射}。
\end{SGAresult}
按照 5.5 的方法，利用 2.1 即可。也可以使用 Serre 的\Term{cohomology}{上同调}准则（EGA II 5.2），无须\Term{descent}{下降}技术便可证明 5.6。

\begin{SGAresult}{推论}{5.7}
在 5.4 的条件下，$f$ 为\Term{integral-morphism}{整态射}（分别为\Term{finite-morphism}{有限态射}、\Term{finite}{有限}且\Term{locally-free}{局部自由}的\Term{morphism}{态射}），当且仅当 $f'$ 有相应性质。
\end{SGAresult}
只需证明充分性，如常可假设 $Y=S$，且 $Y,Y'$ \Term{affine}{仿射}。假设蕴含 $f'$ \Term{affine}{仿射}，由 5.6，$f$ 也\Term{affine}{仿射}，因此 $X$ 以及 $X'$ 都\Term{affine}{仿射}。设 $Y,Y',X,X'$ 的环分别为 $A,A',B,B'=B\otimes_A A'$。有 $B=\varinjlim_i B_i$，其中 $B_i$ 遍历 $B$ 的\Term{finite-type}{有限型} $A$-子代数；于是 $B'=\varinjlim_i B'_i$，其中 $B'_i$ 是 $A'$-代数 $B'$ 的\Term{finite-type}{有限型}子代数。若 $B'$ 在 $A'$ 上整，
\SourcePage{183}
则 $B'_i$ 都是 $A'$ 上的\Term{finite-type}{有限型}模。由于 $A'$ 在 $A$ 上\Term{faithfully-flat}{忠实平坦}，$B_i$ 都是 $A$ 上的\Term{finite-type}{有限型}模，即 $B$ 在 $A$ 上整。同样可见，若 $B'$ 在 $A'$ 上\Term{finite}{有限}，$B$ 就在 $A$ 上\Term{finite}{有限}。对于“\Term{locally-free}{局部自由}\Term{finite-type}{有限型}”，也有同样结论，参见 1.11。

\begin{SGAresult}{推论}{5.8}
在 5.4 的条件下，假设 $f$ \Term{quasi-compact}{拟紧}，$L$ 为 $X$ 上的\Term{invertible-module}{可逆模}，$L'$ 为其在 $X'$ 上的\Term{inverse-image}{逆像}。则 $L$ 相对于 $f$ \Term{ample}{丰沛}（分别为\Term{very-ample}{很丰沛}），当且仅当 $L'$ 相对于 $f'$ 有相应性质。
\end{SGAresult}
只需证明充分性。关于 $L$ 的假设无论如何蕴含 $f'$ \Term{separated}{分离}，故由 4.8，$f$ \Term{separated}{分离}。由于 $f$ \Term{quasi-compact}{拟紧}且 $g:Y'\to Y$ \Term{flat}{平坦}，用\Term{affine}{仿射}\Term{covering}{覆盖}计算\Term{direct-image}{正像}可知，对任意整数 $n$，有\Term{isomorphism}{同构}
\[
g^*\bigl(f_*(L^{\otimes n})\bigr)\xrightarrow{\sim}f'_*(L'^{\otimes n}),
\]
因而有\Term{isomorphism}{同构}
\[
g^*(\mathcal S)\xrightarrow{\sim}\mathcal S',
\]
其中 $\mathcal S$（分别 $\mathcal S'$）表示 $Y$（分别 $Y'$）上的\Term{quasi-coherent}{拟凝聚}\Term{grading}{分次}代数，其\Term{grading}{分次}分量之和为 $n\geq0$ 时各 $f_*(L^{\otimes n})$（分别 $f'_*(L'^{\otimes n})$）的\Term{direct-sum}{直和}。注意，对每个 $n\geq0$，典范\Term{homomorphism}{同态} $f'^*(\mathcal S'_n)\to L'^{\otimes n}$ 的\Term{cokernel}{余核}，是 $f^*(\mathcal S_n)\to L^{\otimes n}$ 的\Term{cokernel}{余核}沿 $X'\to X$ 的\Term{inverse-image}{逆像}，故其\Term{support}{支集} $Z'_n$ 是\Term{support}{支集} $Z_n$ 的\Term{inverse-image}{逆像}。若 $L'$ \Term{ample}{丰沛}，则各 $Z'_n$ 的交为空；由于 $X'\to X$ \Term{surjectivity}{满射}，各 $Z_n$ 的交也为空，即存在典范\Term{morphism}{态射}
\[
j:X\longrightarrow\operatorname{Proj}(\mathcal S)
\]
（EGA II 3）。而类似的\Term{morphism}{态射}
\[
j':X'\longrightarrow\operatorname{Proj}(\mathcal S')
\]
就是前一\Term{morphism}{态射}沿 $Y'\to Y$ 作\Term{base-change}{基变换}所得（同上）。此时，$L'$ 相对于 $f'$ \Term{ample}{丰沛}，就是说 $j'$ 是\Term{immersion}{浸入}；由于 $f'$ \Term{quasi-compact}{拟紧}，此\Term{immersion}{浸入}也必\Term{quasi-compact}{拟紧}。由 5.5，$j$ 为\Term{immersion}{浸入}，即 $L$ 相对于 $f$ \Term{ample}{丰沛}。对于“\Term{very-ample}{很丰沛}”，论证完全类似：只取上面 $n=1$ 的情形，并把 $\operatorname{Proj}(\mathcal S)$ 换成与 $\mathcal S_1$ 相伴的\Term{projective}{射影}丛 $\mathbb P(\mathcal S_1)$。

回顾（EGA II 5.1.1），\Term{quasi-compact}{拟紧}\Term{morphism}{态射} $f$ 称为\Term{quasi-affine-morphism}{拟仿射态射}，是指 $Y$ 中每个\Term{affine-open}{仿射开集} $U$ 的\Term{inverse-image}{逆像} $f^{-1}(U)$ 都\Term{isomorphism}{同构}于某个\Term{affine}{仿射}\Term{scheme}{概形}的\Term{open-subscheme}{开子概形}。该处证明了，这等价于 $\mathcal O_X$ 相对于 $f$ \Term{ample}{丰沛}，也等价于它\Term{very-ample}{很丰沛}。于是 5.8 蕴含：

\SourcePage{184}
\begin{SGAresult}{推论}{5.9}
在 5.4 的条件下，若 $f$ \Term{quasi-compact}{拟紧}，则 $f$ \Term{quasi-affine-morphism}{拟仿射}，当且仅当 $f'$ \Term{quasi-affine-morphism}{拟仿射}。
\end{SGAresult}

\textbf{注记 5.10。}Hironaka 的非\Term{projective-variety}{射影簇}例子说明，可能存在\Term{nonsingular}{非奇异}代数簇之间的\Term{proper-morphism}{固有态射} $f:X\to Y$，其中 $Y$ \Term{projective}{射影}，且 $Y$ 可写成两个开集 $Y_i$ 的并，使 $X_i=X\times_Y Y_i$ 在 $Y_i$ 上\Term{projective}{射影}，但 $f$ 却不是\Term{projective}{射影}\Term{morphism}{态射}。因此，令 $Y'=Y_1\amalg Y_2$，则 $Y'$ 在 $Y$ 上\Term{faithfully-flat}{忠实平坦}且\Term{quasi-compact}{拟紧}，甚至\Term{quasi-finite}{拟有限}；$f':X'\to Y'$ \Term{projective}{射影}，但 $f$ 不\Term{projective}{射影}。所以必须注意：要应用 5.8，从 $f'$ \Term{projective}{射影}推出 $f$ 也\Term{projective}{射影}，必须已经在 $X'$ 上给定一个相对于 $f'$ \Term{ample}{丰沛}的\Term{invertible-module}{可逆模} $L'$，并且它带有相对于 $X'\to X$ 的\Term{descent-datum}{下降数据}。这样才能把 $L'$ 看成 $X$ 上某个\Term{invertible-module}{可逆模} $L$ 的\Term{inverse-image}{逆像}，再由 5.8 得出 $L$ 相对于 $f$ \Term{ample}{丰沛}。不过，当 $g:S'\to S$ \Term{finite}{有限}且\Term{locally-free}{局部自由}时，参见 7.7。

\SGAsection{6}{对有限态射与拟有限态射的应用}
\footnote{对未必\Term{locally-noetherian}{局部诺特}的\Term{prescheme}{预概形}的推广，参见 EGA IV 18.12。}
我们将证明以下两个定理。
\begin{SGAresult}{定理}{6.1}
设 $f:X\to Y$ 为\Term{proper-morphism}{固有态射}，其\Term{fibre}{纤维}\Term{finite}{有限}，且 $Y$ \Term{locally-noetherian}{局部诺特}。则 $f$ 为\Term{finite-morphism}{有限态射}。
\end{SGAresult}
\begin{SGAresult}{定理}{6.2}
设 $f:X\to Y$ \Term{quasi-finite}{拟有限}且\Term{separated-morphism}{分离}，$Y$ \Term{locally-noetherian}{局部诺特}。则 $f$ \Term{quasi-affine-morphism}{拟仿射}，因而特别是\Term{quasi-projective-morphism}{拟射影}的。
\end{SGAresult}
\textbf{注记 6.3。}定理 6.1 是熟知的；在代数簇情形，它由 Chevalley 给出。EGA III 4 中还可找到一个利用“全纯函数定理”的简单证明。这里的证明不用后一定理，而使用\Term{descent}{下降}理论。它可以作为给读者的附赠，因为在证明 6.2 的同时，几乎不费力就能得到它。还需回顾（EGA III 4 或 [1]），由“全纯函数定理”导出的 Zariski“主定理”的整体形式断言：若 $f:X\to Y$ \Term{quasi-finite}{拟有限}且\Term{quasi-projective-morphism}{拟射影}，且 $Y$ 为\Term{noetherian}{诺特}的，则 $X$ 与某个\Term{finite}{有限} $Y$-\Term{prescheme}{预概形} $Z$ 的开子\Term{prescheme}{预概形}为 $Y$-\Term{isomorphism}{同构}。因此，将“主定理”与 6.2 合起来，\Term{representable}{可表}述为：
\begin{SGAresult}{推论}{6.4}
设 $f:X\to Y$ \Term{quasi-finite}{拟有限}且\Term{separated}{分离}，$Y$ 为\Term{noetherian}{诺特}的。则 $X$ 与某个\Term{finite}{有限} $Y$-\Term{prescheme}{预概形} $Z$ 的开子\Term{prescheme}{预概形}为 $Y$-\Term{isomorphism}{同构}。
\end{SGAresult}

\SourcePage{185}
6.2 对\Term{descent}{下降}理论的另一个有趣推论将在 7.9 给出。

\textbf{6.1 与 6.2 的证明。}我们采用以下容易证明的事实。\footnote{参见 EGA IV 8。}
\begin{SGAresult}{引理}{6.5}
设 $X$ 为\Term{locally-noetherian}{局部诺特}\Term{prescheme}{预概形} $Y$ 上的\Term{finite-type}{有限型}\Term{prescheme}{预概形}，$y\in Y$。则存在 $y$ 的\Term{open-neighborhood}{开邻域} $U$，使 $X|_U$ 在 $U$ 上\Term{finite}{有限}（分别\Term{quasi-affine-morphism}{拟仿射}，等等），当且仅当 $X\times_Y\operatorname{Spec}(\mathcal O_y)$ 在 $\operatorname{Spec}(\mathcal O_y)$ 上\Term{finite}{有限}（分别\Term{quasi-affine-morphism}{拟仿射}，等等）。
\end{SGAresult}
另一方面，$f:X\to Y$ \Term{finite}{有限}或\Term{quasi-affine-morphism}{拟仿射}都是 $Y$ 上的局部性质。因此，为证明 6.1、6.2，可归结为 $Y$ 是一个\Term{local-ring}{局部环}的谱的情形；此时特别地，$Y$ 维数\Term{finite}{有限}。注意，\Term{prescheme}{预概形} $Y$ 的维数定义为其各\Term{local-ring}{局部环}的 \Term{krull-dimension}{Krull 维数}的上确界。对
\[
n=\dim(Y)
\]
作归纳；$n<0$ 时断言显然。因此可假设 $n\geq0$，且断言对所有 $n'<n$ 已成立。再次可假设 $Y$ 是一个 $n$ 维\Term{noetherian-local-ring}{诺特局部环} $A$ 的谱。6.1、6.2 的假设在\Term{base-change}{基变换}下保持——开头的\Term{reduced}{约化}已用到了这一点——故在\Term{base-change}{基变换} $\operatorname{Spec}(\widehat A)\to\operatorname{Spec}(A)$ 后仍然成立。后者\Term{faithfully-flat}{忠实平坦}且\Term{quasi-compact}{拟紧}，因此由 5.7、5.9，可进一步归结为 $A$ 完备的情形。利用 $A$ 上每个在 $A$ 上\Term{quasi-finite}{拟有限}的\Term{noetherian-local-ring}{诺特局部环} $B$ 都在 $A$ 上\Term{finite}{有限}，以及 $X$ 在 $Y$ 上\Term{separated}{分离}且 $y$ 处的\Term{fibre}{纤维}由孤立点组成，可得到分解
\[
X=X'\amalg X'',
\]
其中 $X'$ 在 $Y$ 上\Term{finite}{有限}，而 $X''$ 在 $y$ 处的\Term{fibre}{纤维}为空。若 $X$ 在 $Y$ 上\Term{proper}{固有}，则 $X''$ 也如此，故它在 $Y$ 中的像闭；这个像不含 $y$，因此为空。所以 $X''=\varnothing$，$X=X'$，从而 $X$ 在 $Y$ 上\Term{finite}{有限}，6.1 得证。注意，此处并不需要归纳假设。若 $X$ 在 $Y$ 上\Term{quasi-finite}{拟有限}，$X''$ 也如此，而 $X''$ 实际位于 $Y$ 的开子集 $Y-\{y\}$ 上，后者维数小于 $n$。由归纳假设，$X''$ 在 $Y-\{y\}$ 上\Term{quasi-affine-morphism}{拟仿射}，因而在 $Y$ 上也\Term{quasi-affine-morphism}{拟仿射}。$X'$ 显然同样如此，故它们的不交并 $X$ 也如此，6.2 得证。

\textbf{注记 6.6。}不再假设 $Y$ \Term{locally-noetherian}{局部诺特}时，只要明确要求 $f$ \Term{finite-presentation}{有限表示}（参见 3.6），定理 6.1、6.2 仍然成立。事实上，仍可假设 $Y$ \Term{affine}{仿射}，
\SourcePage{186}
此时不难验证，情形 $f:X\to Y$ 是由某个情形 $f_0:X_0\to Y_0$ 沿\Term{base-change}{基变换} $Y\to Y_0$ 得到的，其中 $f_0$ 满足与 $f$ 相同的假设，而 $Y_0$ 为\Term{noetherian}{诺特}的。由 6.1 或 6.2，$f_0$ \Term{finite}{有限}或\Term{quasi-affine-morphism}{拟仿射}，因此 $f$ 也有相应性质。这类论证经常可以用来摆脱\Term{noetherian}{诺特}性假设，而这些假设在应用中总会变得碍事。

\SGAsection{7}{下降数据的有效性准则}
如常考虑\Term{prescheme}{预概形}\Term{morphism}{态射}
\[
g:S'\longrightarrow S
\]
以及一个 $S'$-\Term{prescheme}{预概形} $X'$。根据第 VII 讲第 9 节的一般事实，\footnote{该节并不存在；参见 Bourbaki 讨论班第 212 讲（1962 年）。}在 $X'$ 上给定相对于 $g$ 的\Term{descent-datum}{下降数据}，等价于给定一个\Term{equivalence-pair}{等价态射对} [3]
\[
q_1,q_2:X''\rightrightarrows X',
\]
使结构\Term{morphism}{态射} $X'\to S'$ 与此\Term{morphism}{态射}对及 $g$ 定义的\Term{equivalence-pair}{等价态射对}
\[
p_1,p_2:S''=S'\times_S S'\rightrightarrows S'
\]
相容，并使由相应图
\[
\begin{tikzcd}[column sep=large,row sep=large]
X'\arrow[d]&X''\arrow[l,shift left=.65ex]\arrow[l,shift right=.65ex]\arrow[d]\\
S'&S''\arrow[l,shift left=.65ex]\arrow[l,shift right=.65ex]
\end{tikzcd}
\]
分别取 $p_1,q_1$ 或 $p_2,q_2$ 得到的两个方块都是\Term{cartesian-square}{笛卡尔方块}；由对称性，只要求其中一个也一样。这个\Term{descent-datum}{下降数据}所提出的\Term{descent}{下降}问题的一个解，即一个 $S$ 上的对象 $X$，连同与\Term{descent-datum}{下降数据}相容的\Term{isomorphism}{同构} $X\times_S S'\xleftarrow{\sim}X'$，等价于给定一个\Term{cartesian-square}{笛卡尔方块}
\[
\begin{tikzcd}[column sep=large,row sep=large]
X\arrow[d]&X'\arrow[l,"h"']\arrow[d]\\
S&S'\arrow[l,"g"']
\end{tikzcd}
\]
且满足 $hq_1=hq_2$。

\Term{faithfully-flat}{忠实平坦}且\Term{quasi-compact}{拟紧}的\Term{morphism}{态射}所成的集合在\Term{base-change}{基变换}下保持，而且由 5.3，这样的\Term{morphism}{态射}是\Term{effective-epimorphism}{有效满态射}；因此由一般理论 [D] 得到：

\SourcePage{187}
\begin{SGAresult}{命题}{7.1}
假设 $g:S'\to S$ \Term{faithfully-flat}{忠实平坦}且\Term{quasi-compact}{拟紧}。则 $X'$ 上相对于 $g$ 的\Term{descent-datum}{下降数据}有效，当且仅当它定义的等价关系 $R=(q_1,q_2)$ 有效，即商 $X'/R$ 存在，且 $X''$ 成为 $X'$ 在 $X'/R$ 上的\Term{fibre}{纤维}平方，并且典范\Term{morphism}{态射} $X'\to X'/R$ \Term{faithfully-flat}{忠实平坦}且\Term{quasi-compact}{拟紧}。
\end{SGAresult}
因此，\Term{descent-datum}{下降数据}的有效性问题是等价关系图的有效性问题的一个特例，本节给出的若干有效性准则也可以由此得到。不过，在\Term{descent}{下降}的语境中，我们有定理 2.1，它说明：若 $X'$ 在 $S'$ 上\Term{affine}{仿射}，则 $X'$ 上相对于 $g$ 的每个\Term{descent-datum}{下降数据}都有效。对于一般\Term{flat}{平坦}等价关系图的取商问题，没有与此对应的断言。这里给出的所有有效性准则，也都可以看作从这一准则导出。

设 $U'$ 为 $X'$ 的子\Term{prescheme}{预概形}，或更一般地，为\Term{category}{范畴} $\operatorname{Sch}$ 中 $X'$ 的\Term{subobject}{子对象}。若能在 $U'$ 上赋予相对于 $g$ 的\Term{descent-datum}{下降数据}，使\Term{immersion}{浸入} $U'\to X'$ 与\Term{descent-datum}{下降数据}相容，就称 $U'$ 对 $X'$ 上的\Term{descent-datum}{下降数据}稳定。这也意味着 $U'$ 沿 $q_1,q_2$ 在 $X''$ 中的两个\Term{inverse-image}{逆像}相同；也就是说，$U'$ 对等价关系 $R$ 稳定。显然，此时 $U'$ 上的\Term{descent-datum}{下降数据}唯一，称为由 $X'$ 上的\Term{descent-datum}{下降数据}诱导的\Term{descent-datum}{下降数据}。于是：

\begin{SGAresult}{命题}{7.2}
设 $(X'_i)$ 是 $X'$ 的一个开\Term{covering}{覆盖}，其中每个 $X'_i$ 对\Term{descent-datum}{下降数据}稳定。则 $X'$ 上的\Term{descent-datum}{下降数据}有效，当且仅当各 $X'_i$ 上诱导的\Term{descent-datum}{下降数据}都有效。
\end{SGAresult}
例如，这容易由 7.1 推出；证明细节留给读者。

\begin{SGAresult}{推论}{7.3}
设 $(S_i)$ 为 $S$ 的一个开\Term{covering}{覆盖}，对每个 $i$，令 $S'_i,X'_i$ 分别由 $S',X'$ 沿 $S_i\to S$ 作\Term{base-change}{基变换}得到。则 $X'$ 上的\Term{descent-datum}{下降数据}有效，当且仅当对每个 $i$，$X'_i$ 上相对于 $g_i:S'_i\to S_i$ 的\Term{descent-datum}{下降数据}有效。
\end{SGAresult}
实际使用时，这个准则总能把问题归结为 $S$ \Term{affine}{仿射}的情形。如果 $S'$ 也\Term{affine}{仿射}——这是应用中最常见的情形——则有：
\begin{SGAresult}{推论}{7.4}
假设 $S,S'$ 都\Term{affine}{仿射}。则 $X'$ 上的\Term{descent-datum}{下降数据}有效，当且仅当 $X'$ 是一些对\Term{descent-datum}{下降数据}稳定的\Term{affine-open}{仿射开集} $X'_i$ 的并。
\end{SGAresult}
充分性来自 7.2 及以下事实：若 $X'_i$ \Term{affine}{仿射}，它就在 $S'$ 上\Term{affine}{仿射}，因而可以应用 2.1。对于必要性，注意若 $X'$ 来自 $X$，且 $X$ 被\Term{affine-open}{仿射开集} $X_i$ \Term{covering}{覆盖}，
\SourcePage{188}
则 $X'_i=X_i\times_S S'$ 是对\Term{descent-datum}{下降数据}稳定的\Term{affine-open}{仿射开集}，并\Term{covering}{覆盖} $X'$。

\begin{SGAresult}{推论}{7.5}
设 $g:S'\to S$ 为\Term{faithfully-flat}{忠实平坦}、\Term{quasi-compact}{拟紧}且\Term{radicial}{根性}的\Term{morphism}{态射}。则 $g$ 是\Term{effective-descent-morphism}{有效下降态射}，即对 $S'$ 上的任意 $X'$，其上相对于 $g:S'\to S$ 的每个\Term{descent-datum}{下降数据}都有效。
\end{SGAresult}
事实上，由 7.3，可假设 $S$ \Term{affine}{仿射}。$S'$ 在 $S$ 上\Term{radicial}{根性}，从而\Term{separated}{分离}，因此 $S'$ \Term{separated}{分离}。另一方面，对每个 $x'\in X'$，等价关系 $R$ 所定义的集合意义下的等价关系之等价类
\[
R(x')=q_2\bigl(q_1^{-1}(x')\bigr)
\]
只含一个点：因为 $g$ \Term{radicial}{根性}，沿 $S'\to S$ 作\Term{base-change}{基变换}所得的 $p_1,p_2$ 也\Term{radicial}{根性}，再由它们沿 $X'\to S''$ 作\Term{base-change}{基变换}所得的 $q_1,q_2$ 也\Term{radicial}{根性}。因此，$X'$ 的每个开集都对\Term{descent-datum}{下降数据}稳定。取 $X'$ 的\Term{affine}{仿射}开\Term{covering}{覆盖} $(X'_i)$；由于 $S'$ \Term{separated}{分离}，这些开集在 $S$ 上\Term{affine}{仿射}，因此由 2.1，它们上面诱导的\Term{descent-datum}{下降数据}有效。最后应用 7.2 即得结论。

注意，7.5 给出了\Term{prescheme}{预概形}\Term{category}{范畴}中已知的唯一一种\Term{effective-descent-morphism}{有效下降态射}；即使只考虑\Term{noetherian}{诺特}\Term{scheme}{概形}，或某个域上的\Term{finite-type}{有限型}\Term{scheme}{概形}，这大概也确实是唯一的情形。

若假设 $S$ \Term{locally-noetherian}{局部诺特}且 $S'$ 在 $S$ 上\Term{finite-type}{有限型}，则 7.5 也是以下结论的特例；这个结论推广了 Weil 的 Galois \Term{descent}{下降}与 Cartier 的不\Term{separable}{可分}\Term{descent}{下降}。
\begin{SGAresult}{推论}{7.6}
设 $g:S'\to S$ \Term{finite}{有限}、\Term{locally-free}{局部自由}且为\Term{surjectivity}{满射}，即它由 $S$ 上一个作为模\Term{locally-free}{局部自由}\Term{finite-type}{有限型}的代数定义；因此 $g$ \Term{faithfully-flat}{忠实平坦}且\Term{quasi-compact}{拟紧}，从而是\Term{descent-morphism}{下降态射}。设 $X'$ 为带有\Term{descent-datum}{下降数据}的 $S'$-\Term{prescheme}{预概形}。则此\Term{descent-datum}{下降数据}有效，当且仅当对每个 $x'\in X'$，等价类 $R(x')=q_2(q_1^{-1}(x'))$ 包含在某个\Term{affine-open}{仿射开集}中；若 $X'$ 在 $S'$ 上\Term{quasi-projective-morphism}{拟射影}，这一条件自动满足。
\end{SGAresult}
括号中的说明源自以下事实：若 $s$ 是 $x'$ 下方的 $S$ 中的点，则 $R(x')$ \Term{finite}{有限}且包含在\Term{fibre}{纤维} $X'_s$ 中。另一方面，由于 $X'$ 在 $S'$ 上\Term{quasi-projective-morphism}{拟射影}，且 $S'$ 在 $S$ 上\Term{finite}{有限}，$X'$ 在 $S$ 上\Term{quasi-projective-morphism}{拟射影}，因此 $X'$ 在 $S$ 上的每条\Term{fibre}{纤维}都包含在一个\Term{affine-open}{仿射开集}中。

\Term{affine}{仿射}\Term{scheme}{概形}中的每个\Term{finite}{有限}子集都有一个由\Term{affine}{仿射}邻域组成的基本邻域系，故将问题限制到 $S$ 的某个\Term{affine-open}{仿射开集}上时，不会失去上述假设。由 7.3，便归结为 $S$ \Term{affine}{仿射}的情形。由 7.4，只需证明 $x'$ 位于某个对\Term{descent-datum}{下降数据}稳定的\Term{affine-open}{仿射开集}中。设 $U$ 为包含 $R(x')$ 的\Term{affine-open}{仿射开集}，则饱和集
\[
R(X'-U)=q_2\bigl(q_1^{-1}(X'-U)\bigr)
\]
不与 $R(x')$ 相交。又因 $q_2$ \Term{finite}{有限}——$g$ \Term{finite}{有限}，因而 $p_2$ 也\Term{finite}{有限}——从而为闭映射，右端是 $X'$ 的闭子集。记它在 $X'$ 中的补集为 $U'$，
\SourcePage{189}
则 $U'$ 为饱和开集，且
\[
R(x')\subset U'\subset U.
\]
这里 $U$ \Term{affine}{仿射}，但 $U'$ 还不能先验地断言为\Term{affine}{仿射}。\Term{affine}{仿射}\Term{scheme}{概形} $U$ 中的\Term{finite}{有限}子集 $R(x')$ 有一个由形如 $U_f$ 的\Term{affine}{仿射}邻域组成的基本邻域系。将 $f$ 换成它在 $U'$ 上的限制，可见存在 $\mathcal O_U$ 的\Term{section}{截面} $f$，使
\[
R(x')\subset U'_f,\qquad U'_f\text{ 仿射}.
\]
令 $U''=q_1^{-1}(U')=q_2^{-1}(U')$，仍以 $q_1,q_2$ 表示诱导\Term{morphism}{态射} $U''\to U'$，并考虑
\[
f'=N_{q_2}\bigl(q_1^*(f)\bigr),
\]
其中 $N_{q_2}$ 表示相对于\Term{finite}{有限}\Term{locally-free}{局部自由}\Term{morphism}{态射} $q_2:U''\to U'$ 的\Term{norm}{范数}。\Term{norm}{范数}的构造与\Term{base-change}{基变换}相容，容易由此推出 $f'$ 是不变\Term{section}{截面}：
\[
q_1^*(f')=q_2^*(f').
\]
因此 $U'_{f'}$ 为 $U'$ 的饱和开集。更精确地，记\Term{section}{截面} $f'$ 的零点集为 $Z(f')$，由\Term{norm}{范数}的性质得到
\[
\begin{aligned}
Z(f')&=q_2\bigl(Z(q_1^*(f))\bigr)\\
&=q_2\bigl(q_1^*(Z(f))\bigr)=R(U'-U'_f).
\end{aligned}
\]
所以 $U'_{f'}=U'-Z(f')$ 饱和，包含 $R(x')$，且包含于 $U'_f$。后者\Term{affine}{仿射}，故 $U'_{f'}$ 也\Term{affine}{仿射}，因为它等于 $(U'_f)_{f''}$，其中 $f''=f'|_{U'_{f'}}$。于是得到了一个包含 $R(x')$、从而包含 $x'$ 的饱和\Term{affine-open}{仿射开集}，证明完毕。

注意，上述论证适用于\Term{prescheme}{预概形} $X'$ 中的任意\Term{finite}{有限}\Term{locally-free}{局部自由}等价关系，甚至只要求为预等价关系，参见 [3]。事实上，7.6 也是\Term{finite}{有限}\Term{locally-free}{局部自由}预等价关系的相应结果的特例，参见该文；下面的 7.7 也有同样的说明。

也可以在得到包含 $x'$ 的饱和\Term{quasi-affine-morphism}{拟仿射}开集 $U'$ 之后，应用 7.9 与 7.2，以免使用\Term{norm}{范数}。

另外，在 7.6 的条件下，若 $X'$ 上的\Term{descent-datum}{下降数据}有效，且 $X'$ 来自 $S$ 上的 $X$，则\Term{morphism}{态射} $X'\to X$ \Term{finite}{有限}、\Term{locally-free}{局部自由}且为\Term{surjectivity}{满射}，因为它由 $g$ 沿 $X\to S$ 作\Term{base-change}{基变换}得到。由此可知（EGA II 6.6.4），若 $X'$ 在 $S'$ 上\Term{quasi-projective-morphism}{拟射影}，从而在 $S$ 上\Term{quasi-projective-morphism}{拟射影}，则 $X$ 在 $S$ 上\Term{quasi-projective-morphism}{拟射影}。$X$ 上的相对\Term{ample}{丰沛}\Term{invertible-sheaf}{可逆层}，可以对 $X'$ 上一个相对于 $S$ \Term{ample}{丰沛}的\Term{invertible-sheaf}{可逆层}取\Term{norm}{范数}而得到；这里“相对于 $S$”也可换成“相对于 $S'$”，两者等价。于是得到：

\SourcePage{190}
\begin{SGAresult}{推论}{7.7}
\Term{finite}{有限}、\Term{locally-free}{局部自由}且\Term{surjectivity}{满射}的\Term{morphism}{态射} $g:S'\to S$，对于由在其他\Term{prescheme}{预概形}上\Term{quasi-projective-morphism}{拟射影}的\Term{prescheme}{预概形}组成的\Term{fibred-category}{纤维范畴}，是\Term{effective-descent-morphism}{有效下降态射}。也就是说，对每个在 $S'$ 上\Term{quasi-projective-morphism}{拟射影}的 $X'$，它上面相对于 $g$ 的每个\Term{descent-datum}{下降数据}都有效，而且\Term{descent}{下降}所得的 $S$-\Term{prescheme}{预概形} $X$ 在 $S$ 上\Term{quasi-projective-morphism}{拟射影}。
\end{SGAresult}

\begin{SGAresult}{命题}{7.8}
设 $g:S'\to S$ \Term{faithfully-flat}{忠实平坦}且\Term{quasi-compact}{拟紧}。考虑由在\Term{prescheme}{预概形} $T$ 上\Term{quasi-compact}{拟紧}的\Term{prescheme}{预概形} $Z$，连同其上一个相对于 $T$ \Term{ample}{丰沛}的\Term{invertible-sheaf}{可逆层}组成的\Term{fibred-category}{纤维范畴}，则 $g$ 是这个\Term{category}{范畴}的\Term{effective-descent-morphism}{有效下降态射}。特别地，设 $X'$ 为 $S'$ 上的\Term{prescheme}{预概形}，并带有相对于 $g:S'\to S$ 的\Term{descent-datum}{下降数据}；$L'$ 为 $X'$ 上相对于 $S'$ \Term{ample}{丰沛}的\Term{invertible-sheaf}{可逆层}，且也带有相对于 $X'$ 上所给\Term{descent-datum}{下降数据}的\Term{descent-datum}{下降数据}，即给定一个 $q_1^*(L')$ 到 $q_2^*(L')$ 的\Term{isomorphism}{同构}，满足通常的传递性条件。则 $X'$ 上的\Term{descent-datum}{下降数据}有效，而且 $L'$ \Term{descent}{下降}到所得\Term{prescheme}{预概形} $X$ 上的\Term{invertible-sheaf}{可逆层} $L$ 相对于 $S$ \Term{ample}{丰沛}。
\end{SGAresult}
证明与 5.8 完全类似。由 $L'$ 定义的 $S'$ 上的\Term{quasi-coherent}{拟凝聚}\Term{grading}{分次}代数 $\mathcal S'$ 也带有\Term{descent-datum}{下降数据}，因而可利用 1.1 构造 $S$ 上的\Term{quasi-coherent}{拟凝聚}\Term{grading}{分次}代数 $\mathcal S$。于是得到 $S$ 上的 $P=\operatorname{Proj}(\mathcal S)$，使 $P'=\operatorname{Proj}(\mathcal S')$ 连同其\Term{descent-datum}{下降数据}等同于 $P\times_S S'$。按假设，$X'$ 等同于 $P'$ 的一个开子集，它必对 $P'$ 上的\Term{descent-datum}{下降数据}稳定，故 $X'$ 上的\Term{descent-datum}{下降数据}也有效，所得\Term{descent}{下降}\Term{prescheme}{预概形}是 $P$ 的一个开子集。细节留给读者。特别令 $L'=\mathcal O_{X'}$，得到：

\begin{SGAresult}{推论}{7.9}
设 $g:S'\to S$ \Term{faithfully-flat}{忠实平坦}且\Term{quasi-compact}{拟紧}，$X'$ 为 $S'$ 上的\Term{quasi-affine-morphism}{拟仿射}\Term{prescheme}{预概形}。则 $X'$ 上相对于 $g$ 的每个\Term{descent-datum}{下降数据}都有效，而且\Term{descent}{下降}所得的\Term{prescheme}{预概形} $X$ 在 $S$ 上\Term{quasi-affine-morphism}{拟仿射}。
\end{SGAresult}
由 6.2，特别当 $S'$ \Term{locally-noetherian}{局部诺特}，且 $X'$ 在 $S'$ 上\Term{quasi-finite}{拟有限}且\Term{separated}{分离}时，可以应用这一结果。更一般地，$S'$ 任意，而 $X'$ 在 $S'$ 上\Term{finite-presentation}{有限表示}、\Term{quasi-finite}{拟有限}且\Term{separated}{分离}时，也可应用，参见 6.6。

\textbf{注记 7.10。}本节给出的结果囊括了目前已知的有效性准则，甚至大概也囊括了所有有用的有效性准则。\footnote{这一看法后来部分被证明有误，例如参见 J. P. Murre，Bourbaki 讨论班第 294 讲（附录），1965 年 5 月，以及若干特殊结果，特别是 Néron 与 Raynaud 的结果。\Term{group-scheme}{群概形}的\Term{descent}{下降}参见 M. Raynaud 的博士论文（1968 年）。}以下反例可以支持这一断言：

\SourcePage{191}
\begin{enumerate}[label=(\roman*)]
\item 若 $S$ 是一个域的谱，$S'$ 是其二次 Galois \Term{group-extension}{扩张}的谱，则可找到 $S'$ 上的 $X'$，使它在 $S'$ 上\Term{proper}{固有}且\Term{smooth}{光滑}，维数为 $3$，并带有一个无效的\Term{descent-datum}{下降数据}（Serre）。
\item 可以找到 $S$，它是一个三维\Term{regular-local-ring}{正则局部环}的谱；如有需要，可取为给定特征的域上某个代数\Term{scheme}{概形}的\Term{local-ring}{局部环}。还可找到 $S$ 的一个群为 $\mathbb Z/2\mathbb Z$ 的\Term{principal-covering}{主覆盖} $T$，使得若以 $t$ 表示 $T$ 中位于 $S$ 的\Term{closed-point}{闭点} $s$ 上方的一个点，并令 $S'=T-s$，则存在 $S'$ 上的一个\Term{projective}{射影}且\Term{regular}{正则}的 $X'$，它带有相对于 $g:S'\to S$ 的\Term{descent-datum}{下降数据}，而这个\Term{descent-datum}{下降数据}无效。
\end{enumerate}
这些构造利用 Hironaka 的非\Term{projective-variety}{射影簇}例子。对于 (i)，只需利用如下事实：在 $k$ 上，可以找到一个三维\Term{proper}{固有}\Term{smooth}{光滑}\Term{scheme}{概形} $X_0$，群 $G=\mathbb Z/2\mathbb Z$ 在其上无惯性地作用，而且存在 $X_0$ 的两个在 $S$ 上有理的点 $a,b$，它们在 $G$ 作用下相合，却不同时包含在任何\Term{affine-open}{仿射开集}中。令 $X'=X_0\times_k k'$，利用 $G$ 对两个因子的作用，使其作用于 $X'$，便得到 $X'$ 上相对于 $g:\operatorname{Spec}(k')\to\operatorname{Spec}(k)$ 的\Term{descent-datum}{下降数据}。在 $a$ 和 $b$ 上方分别恰有一个点 $a'$ 和 $b'$，其\Term{residue-field}{剩余域}\Term{group-extension}{扩张}为二次\Term{group-extension}{扩张}。由于 $X'\to X_0$ 与 $G$ 的作用相容，$a',b'$ 在 $G$ 作用下相合。它们不能同时包含在一个\Term{affine-open}{仿射开集} $U'$ 中，否则
\[
U=X_0-\operatorname{Im}(X'-U')
\]
就是 $X_0$ 中包含 $(a,b)$ 的开集，其在 $X'$ 中的\Term{inverse-image}{逆像}包含于 $U'$，从而\Term{quasi-affine-morphism}{拟仿射}；于是 $U$ \Term{quasi-affine-morphism}{拟仿射}，故 $(a,b)$ 在 $U$ 中会有一个\Term{affine}{仿射}邻域。

对于 (ii)，利用 Hironaka 的例子中以下事实：$X_0$ 是某个 $k$-\Term{projective}{射影}\Term{scheme}{概形} $Y$ 上的\Term{proper}{固有}\Term{prescheme}{预概形}，且 $Y$ 在 $k$ 上\Term{smooth}{光滑}。\Term{morphism}{态射} $f:X_0\to Y$ 还是双有理的，不过这一点无关紧要。群 $G$ 也作用于 $Y$，并与它在 $X_0$ 上的作用相容。最后，令 $S'=Y-f(b)$，$X'=X_0|_{S'}$，则 $X'$ 在 $S'$ 上\Term{projective}{射影}。由于 $G$ 在 $X_0$ 上的作用与它在 $Y$ 上的作用相容，$X_0$ 自然带有相对于典范\Term{morphism}{态射} $Y\to S=Y/G$ 的\Term{descent-datum}{下降数据}。这个\Term{descent-datum}{下降数据}无效，因为 $(a,b)$ 不同时包含在一个\Term{affine-open}{仿射开集}中。由它在 $X'$ 上诱导的相对于 $g:S'\to S$ 的\Term{descent-datum}{下降数据}也无效，这一点容易验证。

\SGAsection{8}{参考文献}
\begin{description}[leftmargin=2.8em,labelsep=.5em]
\item[{[D]}] J. Giraud，\textit{Méthode de la descente}（《\Term{descent}{下降}方法》），法国数学会专刊第 2 号，1964 年。
\item[{[1]}] A. Grothendieck，Bourbaki 讨论班：\textit{Géométrie formelle et Géométrie algébrique}（《\Term{formal-geometry}{形式几何}与代数几何》），1959 年 5 月，第 182 讲。
\item[{[2]}] A. Grothendieck，Bourbaki 讨论班：\textit{Technique de descente et Théorèmes d'existence I}（《\Term{descent}{下降}技术与存在定理 I》），1959 年 12 月，第 190 讲。
\SourcePage{192}
\item[{[3]}] A. Grothendieck，Bourbaki 讨论班：\textit{Technique de descente et Théorèmes d'existence III}（《\Term{descent}{下降}技术与存在定理 III》），1961 年 2 月，第 212 讲。
\end{description}
